\section{引言}

分钟级交易数据使得短时市场状态的连续监测成为可能。对于 A 股市场而言，单个证券的短时成交承接、价格振幅和收益波动往往会在市场压力时段同步恶化，进而形成跨证券的共同压力。本文关注的问题是：仅使用 1 分钟 OHLCV 与成交额等可得高频信息，是否能够对未来 5 分钟和 10 分钟的短时流动性压力形成可验证的样本外预警。

本文基于 1 分钟 OHLCV 与成交额数据构造短时流动性压力代理指标，重点刻画市场成交、价格振幅、局部波动和成交承接状态共同反映的短时压力变化。该定位决定了本文的经验结论应被理解为“市场状态信息对短时压力事件的预警能力”。围绕这一定位，本文形成三层证据链：首先构造 LSI 与 MarketLSI，定义可验证的未来压力标签；其次使用 RGARCH-CARR-SK 和 QVAR 分别刻画动态风险路径与尾部分位传导；最后用 SMARTBoost 检验样本外预警效果。

\begin{table}[tbp]
\centering
\caption{研究框架、技术实现与解释边界}
\label{tab:model-framework}
\small
\begin{threeparttable}
\begin{tabularx}{\textwidth}{p{0.16\textwidth}p{0.26\textwidth}Y}
\toprule
模块 & 技术实现 & 本文用途与解释边界 \\
\midrule
压力指数构造 & LSI 与 MarketLSI & 将分钟级 OHLCV 与成交额信息压缩为短时压力代理指标；不等同于订单簿真实流动性。 \\
动态风险刻画 & \mbox{RGARCH-CARR-SK} & 刻画 MarketLSI 压力创新的条件风险和高阶矩路径；用于风险度量而非收益率预测。 \\
尾部分位传导 & QVAR & 描述市场状态变量在不同分位点下的联动和情景响应；情景模拟不作严格因果识别。 \\
样本外预警 & SMARTBoost & 检验无泄漏历史特征对未来 H5/H10 压力事件的预测能力；作为唯一正式机器学习预警模型。 \\
\bottomrule
\end{tabularx}
\begin{tablenotes}[flushleft]
\footnotesize
\item 注：三类模型形成递进证据链，而不是并列的模型竞赛。
\end{tablenotes}
\end{threeparttable}
\end{table}


表 \ref{tab:model-framework} 概括了本文方法体系的功能分工。RGARCH-CARR-SK 与 QVAR 主要服务于风险刻画和机制描述，SMARTBoost 则承担正式的样本外预警检验。这样的安排避免将三类模型简单处理为预测比赛，也使本文能够同时回答“压力如何刻画”“压力如何传导”和“压力能否预警”三个层次的问题。

本文的主要贡献体现在三方面。第一，基于分钟级 OHLCV 与成交额构造了面向短时压力事件的 LSI 指标，并严格使用训练期历史分布确定标准化参数与标签阈值。第二，将 RGARCH-CARR-SK 的 realized measure 驱动风险递推结构迁移到 MarketLSI 压力风险序列，形成适合本文数据对象的动态风险度量。第三，在排除 CrossStress 等未来标签聚合变量后，使用 SMARTBoost 对 \(\StressHfive\) 与 \(\StressHten\) 进行时间滚动样本外预警，重点报告 PR-AUC、Top 5\% 高风险组命中率与 lift 等更适合稀有压力事件的指标。
